\section{从一个简单循环到社区平台}

Reddit 的故事适合用来观察互联网系统怎样从产品原型长成社会基础设施。最初的目标并不是“构建数万个社区”，而是让用户发现新的网络内容：提交链接、投票排序、围绕链接讨论。这个循环足够简单，创始人可以在数周内实现；它又足够开放，用户能逐渐把新闻聚合器改造成身份、规范和知识都不同的社区网络。理解这条演化路径，比背诵创办年份更能解释后续治理与基础设施为什么复杂。

\subsection{Slashdot、Delicious 与最小产品闭环}

讲者把原始产品描述为 Slashdot 和 Delicious 的结合。Slashdot 的价值来自新闻条目下的技术社区讨论，Delicious 则展示用户公开收藏、打标签和共同排序链接的可能。Reddit 去掉传统编辑，把内容选择交给提交与投票，使“谁决定首页”从中心化岗位变成一套可计算的参与规则。此时 subreddit 还不是显式产品主角，社区先在重复互动中出现。

\begin{figure}[H]
\centering
\includegraphics[width=0.94\textwidth]{images/01-origin-product-loop.png}
\caption{Reddit 最初的产品闭环：提交、投票、排序与讨论互相推动。}
\figsource{依据访谈字幕 00:00--05:00 重绘，`images/01-origin-product-loop.png`。}
\end{figure}

\begin{knowledgebox}{读图：循环中的四种状态}
Submit 产生候选内容，vote 聚合用户偏好，rank 决定有限注意力如何分配，discuss 则把一次点击变成持续关系。前三步可以由数据库和排序逻辑完成，第四步会形成规范、声誉和冲突。系统扩张后，难题不再只是把链接排对顺序，而是维护成千上万个讨论空间之间的边界。
\end{knowledgebox}

\begin{importantbox}{术语消化：UGC 与 subreddit}
UGC 是 user-generated content，即用户生成内容，包含帖子、评论、图片、投票和社区规则。Subreddit 是 Reddit 内围绕主题形成的社区单元，拥有自己的版主、规则和成员文化。平台提供全局基础设施，subreddit 负责大量本地治理，两层权力会在边界案例中相互影响。
\end{importantbox}

\subsection{Cold start：创始人何时可以退出循环}

\term{Cold start} 指系统尚无足够用户、内容或互动，导致新用户进入后看不到价值。Reddit 早期首页只展示近 24 小时链接，如果创始人不持续填充，页面可能完全空白。Huffman 和 Alexis Ohanian 因此用脚本抓取链接，再用多个账号提交，让访问者看见一个已经运转的产品。这个做法解决的是展示与学习问题，却也带来真实性边界：种子内容应尽快让位给真实参与，而不能永久伪造社区规模。

\begin{knowledgebox}{课堂提示：一次睡过头成为增长判据}
Huffman 回忆，某个周末他没有执行日常填充脚本，打开网站时却发现首页已经被陌生用户名提交的链接填满。这个瞬间比注册数更有意义：系统在没有创始人当日劳动的情况下继续运转。对双边市场、社区和数据产品而言，真正的 product-market fit 往往表现为关键循环能自行续航。
\end{knowledgebox}

\begin{table}[H]
\centering
\begin{tabular}{L{0.2\textwidth}L{0.31\textwidth}L{0.39\textwidth}}
\toprule
阶段 & 创始人动作 & 应观察的退出信号 \\
\midrule
空白期 & 手工或脚本准备高质量种子 & 新用户能理解产品用途 \\
引导期 & 建立多个主题与互动示例 & 用户开始模仿核心行为 \\
临界期 & 降低人工填充频率 & 首页或供给在无人值守时仍更新 \\
自治期 & 转向反滥用、工具和社区支持 & 留存来自用户之间的价值交换 \\
\bottomrule
\end{tabular}
\caption{冷启动不是一次增长活动，而是创始人劳动逐步退出核心循环的过程。}
\end{table}

\subsection{本章小结}

Reddit 的早期优势不是复杂功能，而是一条用户可接管的循环。创始人可以暂时提供种子供给，但最终判据是未知用户是否持续提交、投票和讨论。循环一旦形成，平台就开始积累社区状态；这些状态会在公司所有权和组织人员多次变化时继续存在。

